%!TEX root = ../main.tex
% =============================================================
%  CAPÍTULO 8 — IMPLICANCIAS, CONCLUSIONES Y RECOMENDACIONES
%                                       (Entrega 8 · 5 % · 20/10)
% =============================================================

\chapter{Implicancias, conclusiones y recomendaciones}\label{cap:conclusiones}

Este capítulo presenta las implicancias del trabajo, sus conclusiones
(respuesta a la pregunta de investigación y a la hipótesis, grado de
cumplimiento de los objetivos, interpretación de los resultados y
limitaciones), los aportes de la tesis, las recomendaciones según el
destinatario y las futuras líneas de investigación.

\section{Implicancias}

\subsection{Innovación y tecnología}

La arquitectura multi-tenant demostró que un mismo sistema puede servir a
tres tipos de usuario con reglas de acceso distintas sin duplicar la base
de datos ni el código. El motor de recomendación y el simulador se
construyeron como módulos propios, separados de la lógica de negocio y
apoyados en un modelo de datos común, sin exigir grandes volúmenes de
datos ni infraestructura costosa; esa separación deja abierta la
posibilidad de reemplazar las reglas ponderadas por un modelo entrenado
cuando existan datos de uso, sin reescribir el resto del sistema.
Respecto de los procesos existentes, la solicitud de contratación dentro
de la plataforma reemplaza el intercambio por WhatsApp, y la
autenticación con roles reemplaza al mecanismo provisorio del prototipo
original. Partir de un prototipo real de gestión, el de Rodríguez
Catering, permitió construir sobre procesos que ya se usaban en lugar de
diseñar el sistema en abstracto.

\subsection{Relevancia social y económica}

\textbf{Social.} Para quien organiza un evento, concentrar la búsqueda, la
comparación y la contratación en un solo lugar reduce el esfuerzo de
coordinar varios canales y la dependencia de intermediarios. Para los
salones y proveedores pequeños, significa acceder a herramientas de
gestión y a visibilidad comercial que habitualmente solo están al
alcance de empresas más grandes, sin una inversión propia en marketing.
Reunir precios y disponibilidad en un solo lugar también reduce la
asimetría de información de quien contrata por primera vez un servicio
que no conoce.

\textbf{Económica.} El modelo combina una suscripción para salones y
proveedores con una comisión del 5\,\% por contratación. En el escenario
base se obtiene un valor actual neto de aproximadamente USD~77.800, una
tasa interna de retorno cercana al 122\,\% y un período de repago de
alrededor de 2,1 años; en el escenario pesimista, el valor actual neto
sigue siendo positivo (USD~33.700). Estos valores son estimaciones
preliminares que dependen de supuestos de adopción y de la tasa de
descuento, y deben leerse como evidencia de viabilidad potencial, no
como una proyección garantizada. El bajo monto de inversión inicial se
explica en gran parte porque el desarrollo es trabajo propio del
tesista; si se incorporara un costo de oportunidad por ese tiempo, los
indicadores bajarían sin modificar la conclusión de viabilidad. La
sostenibilidad depende además de alcanzar una masa crítica de salones y
proveedores que haga atractivo el marketplace para los clientes, y
viceversa, el problema típico de las plataformas de dos lados.

\section{Conclusiones}

\subsection{Hallazgos clave y respuesta a la pregunta de investigación}

La pregunta de investigación indagaba si resulta viable, técnica y
económicamente, integrar en una misma plataforma un marketplace de
servicios, un sistema de recomendación basado en reglas ponderadas y un
Gemelo Digital del evento, y si esa integración despierta interés en
clientes, salones y proveedores, frente al uso de herramientas
fragmentadas. La evidencia reunida permite responder en forma
afirmativa: se verificó la factibilidad técnica y económica de la
integración y el interés de los usuarios relevados. No se midió, en
cambio, la mejora de eficiencia en el uso real, que queda como trabajo
futuro.

En cuanto a la hipótesis, que sostenía que esa integración hace más
simple y eficiente la organización de eventos y que el Gemelo Digital
permite anticipar situaciones y evaluar el impacto de las decisiones, la
respuesta es \textbf{afirmativa, con matices}. Lo que sí se pudo verificar
es lo siguiente:

\begin{itemize}
    \item \textbf{Es técnicamente viable integrar las tres capas.} El
          marketplace, el motor de recomendación y el Gemelo Digital
          funcionan en una misma plataforma, sobre una base de datos
          común y con separación de datos entre salones, proveedores y
          clientes. La suite de pruebas automatizadas (410 pruebas, de las
          cuales 409 pasaron en la última ejecución) verifica, entre otras cosas, que un cliente no puede
          operar sobre el evento de otro, que un salón ocupado en la
          fecha no se recomienda y que la simulación detecta tanto el
          exceso de capacidad como el conflicto de fechas.
    \item \textbf{Existe una necesidad real y un interés declarado.} Según
          la encuesta, quienes organizan un evento hoy recurren a varios
          canales a la vez (redes sociales, recomendaciones de conocidos y
          WhatsApp), el 46,7\,\% tarda entre una y cuatro semanas en buscar
          y comparar proveedores y un 28,9\,\% más de un mes, y el 91,1\,\%
          usaría o probablemente usaría una plataforma que reúna todo. El
          93,3\,\% valoró útil la recomendación asistida y el 97,8\,\% la
          simulación del evento antes de confirmarlo.
    \item \textbf{El modelo es económicamente viable bajo los supuestos
          planteados.} En los tres escenarios analizados el valor actual
          neto es positivo y la tasa interna de retorno supera la tasa de
          corte del 20\,\%, incluso con un 60\,\% de la adopción prevista.
\end{itemize}

El matiz es que \textbf{la mejora de eficiencia frente al uso de
herramientas fragmentadas no se midió de forma directa}. No se realizó un
experimento comparativo (por ejemplo, el tiempo que le lleva a una
persona organizar un mismo evento con la plataforma y sin ella), por lo
que la hipótesis queda respaldada por la viabilidad técnica, la
necesidad relevada y el interés declarado, pero no confirmada con
mediciones de uso real. Tampoco se contrastó el aporte del Gemelo Digital
con eventos reales ya realizados, por lo que su capacidad de anticipar
problemas debe leerse como preliminar.

\subsection{Cumplimiento de los objetivos}

La tabla siguiente retoma cada objetivo específico y resume en qué
medida se alcanzó y con qué evidencia.

\small
\begin{xltabular}{\textwidth}{p{3.1cm} X p{2.1cm} X}
    \caption{Grado de cumplimiento de los objetivos específicos}
    \label{tab:cumplimiento} \\
    \toprule
    \tb{Objetivo} & \tb{Evidencia} & \tb{Grado} & \tb{Observación} \\
    \midrule
    \endfirsthead
    \toprule
    \tb{Objetivo} & \tb{Evidencia} & \tb{Grado} & \tb{Observación} \\
    \midrule
    \endhead
    \bottomrule
    \endfoot
    OE-1: buscar, comparar y contratar desde un mismo lugar & RF-02 y RF-03; pruebas del módulo de marketplace & Alcanzado, con matices & Búsqueda pública, ficha y solicitud de contratación funcionan; el cobro de seña y saldo funciona con Mercado Pago solo en modo de prueba \\
    OE-2: herramientas de gestión para salones y proveedores & RF-06; módulos de clientes, presupuestos y contratos heredados del prototipo de Rodríguez Catering; pruebas de aislamiento entre negocios y de pagos & Alcanzado, con matices & Cada tenant administra solo sus reservas; la gestión de cobros está integrada a Mercado Pago solo en modo de prueba, no en producción \\
    OE-3: gestionar la disponibilidad de servicios y recursos & RF-02 y RF-04 (filtro por disponibilidad); detección de conflicto de fechas en el Gemelo Digital & Alcanzado & La disponibilidad se evalúa por fecha \\
    OE-4: herramientas de recomendación & RF-04 y RNF-02; 8 pruebas del módulo de recomendación & Alcanzado, con matices & Es un conjunto de reglas ponderadas con motivo explicable, no un modelo entrenado; no se validó con usuarios reales \\
    OE-5: Gemelo Digital para analizar escenarios & RF-05; 9 pruebas del módulo de Gemelo Digital & Alcanzado & Simula invitados, fecha y presupuesto sobre el evento propio, sin alterarlo, y permite aplicar el escenario elegido \\
    OE-6: información e indicadores para decidir & Motivo de cada recomendación; costo estimado y avisos de la simulación & Parcial & Hay información para decidir a nivel de evento (motivo de la recomendación, costo estimado y avisos); falta evaluar si alcanza para salones y proveedores \\
\end{xltabular}
\normalsize

De los seis objetivos, cinco se alcanzaron (tres de ellos con matices) y
uno se alcanzó en forma parcial. En ningún caso se trata de
funcionalidades ausentes: las diferencias están en lo que depende de
terceros (el paso de los pagos del modo de prueba a producción), de
datos de uso que todavía no existen (la validación de la recomendación)
o de lo que falta evaluar con usuarios reales (los indicadores para
salones y proveedores).

\subsection{Interpretación de los resultados}

\textbf{Significado en el contexto del estudio.} Solo el 2,2\,\% de
quienes respondieron la encuesta usó un marketplace existente, aunque la
búsqueda de proveedores les lleva semanas y varios canales. Esto sugiere
que la oportunidad no está en la falta de plataformas, sino en la baja
adopción de las existentes y en el esfuerzo de comparar y coordinar. El
interés declarado en la recomendación y en la simulación (93,3\,\% y
97,8\,\%) superó al interés en la plataforma unificada (91,1\,\%), lo que
sugiere que el componente inteligente forma parte de lo que motivaría la
adopción.

\subsection{Limitaciones}

\subsubsection{Técnicas}

El motor de recomendación se basa en reglas ponderadas y no aprende de
los datos, por lo que su calidad depende de que los pesos elegidos sean
razonables. Las pruebas automatizadas verifican el comportamiento
esperado pero no garantizan haber previsto todos los casos. No se midió el
rendimiento con carga real ni se hicieron pruebas de concurrencia,
porque el sistema todavía no está desplegado. Los pagos solo se
probaron en el modo de prueba (sandbox) de Mercado Pago, no con cobros
reales.

\subsubsection{De datos}

La plataforma no cuenta todavía con historial de uso real: las reseñas,
las reservas y las preferencias que alimentarían una recomendación más
precisa recién existirán una vez que haya usuarios. El prototipo de
gestión se construyó a partir de un único caso real, el de Rodríguez
Catering, por lo que no es posible generalizar su desempeño a otros
tipos de proveedor. Las estimaciones de adopción del análisis económico
provienen de fuentes secundarias y de supuestos propios, y no de datos
observados.

\subsubsection{De validación}

La encuesta se distribuyó por conveniencia, a través de redes sociales y
contactos del autor, por lo que sus resultados se interpretan como una
validación exploratoria de la dirección del proyecto y no como una
medición representativa de la población. Se realizaron las tres
entrevistas previstas a especialistas, aplicadas por escrito y no en
vivo. No se hizo una prueba de usabilidad guiada con usuarios
reales sobre el prototipo en funcionamiento, y los resultados fueron
evaluados por el propio desarrollador.

\subsubsection{De mercado}

Los resultados corresponden al mercado argentino, con foco en Buenos
Aires y el Gran Buenos Aires; no pueden extrapolarse sin más a otras
regiones. El análisis económico asume un contexto de inflación y tipo de
cambio que puede variar, y no se evaluó la reacción de las plataformas ya
existentes ni la disposición real de los salones y proveedores a pagar
la suscripción.

\section{Aportes de la tesis}

\subsection{Contribuciones originales}

Dentro de las plataformas relevadas, ninguna combina un marketplace de
servicios para eventos, herramientas de gestión para salones y
proveedores, un motor de recomendación con motivo explicable y un
simulador de escenarios del evento. El trabajo aporta un diseño concreto
y verificable de esa combinación, con una arquitectura multi-tenant que
mantiene separados los datos de cada actor, y una implementación de
referencia con sus pruebas automatizadas. Su contribución no es un
resultado experimental concluyente, sino un punto de partida verificable:
delimita qué es técnicamente viable, qué supuestos económicos sostienen el
modelo y qué datos habría que reunir para validarlo.

\subsection{Beneficios para el campo de estudio}

Para la ingeniería de sistemas aplicada a la gestión de eventos, el
trabajo deja una referencia de arquitectura reutilizable, un análisis de
factibilidad con tres escenarios que puede adaptarse a otros marketplaces
de servicios, y un relevamiento de necesidades de clientes, salones y
proveedores del mercado local, con su instrumento (cuestionario y guías de
entrevista) disponible en los apéndices.

\subsection{Impacto en la comunidad académica}

A la fecha no hay publicaciones ni participaciones en conferencias
derivadas del trabajo. El documento y, una vez definido el repositorio
oficial, el código quedan como material de consulta para futuros
proyectos de la cátedra.

\subsection{Aplicaciones futuras}

El motor de reglas queda como línea base contra la cual comparar un modelo
entrenado, y el instrumento de la encuesta puede reutilizarse con una
muestra más amplia. No se identificaron, en cambio, implicancias para
políticas o regulaciones del área más allá de la protección de datos
personales, tratada en las recomendaciones.

\section{Recomendaciones}

\subsection{Para la práctica profesional}

\begin{enumerate}
    \item \textbf{Desplegar un piloto acotado} con un número pequeño de
          salones y proveedores reales antes de pensar en escala. Es la
          única forma de obtener los datos de uso que hoy faltan y de
          medir el rendimiento bajo carga real.
    \item \textbf{Medir la eficiencia frente a la alternativa fragmentada}
          con un estudio comparativo simple: tiempo y cantidad de
          contactos necesarios para armar un mismo evento con y sin la
          plataforma. Es la medición que permitiría responder del todo la
          hipótesis.
    \item \textbf{Realizar una prueba de usabilidad guiada} con entre tres
          y cinco usuarios de cada tipo (cliente, salón, proveedor) sobre
          el prototipo.
    \item \textbf{Pasar los pagos del modo de prueba a producción}
          recién cuando el piloto lo justifique, y reforzar seguridad,
          facturación y escalabilidad antes de un lanzamiento comercial.
    \item \textbf{Decidir de antemano quién mantiene el sistema} después
          de la defensa, y dejar el repositorio y la documentación en
          condiciones de ser retomados por un tercero.
\end{enumerate}

\subsection{Para investigadores futuros}

Conviene repetir la encuesta con un muestreo probabilístico y en otras
regiones, y diseñar desde el inicio un estudio comparativo entre la
plataforma y las herramientas fragmentadas, con indicadores definidos de
antemano. Si se incorpora un modelo de recomendación entrenado, debe
evaluarse con métricas fijadas antes de ver los resultados y contra la
línea base de reglas ya implementada.

\subsection{Para instituciones y políticas públicas}

La plataforma administra datos personales de clientes, salones y
proveedores, por lo que, antes de operar con usuarios reales, conviene
revisar con asesoramiento legal las obligaciones que impone la Ley 25.326
de Protección de Datos Personales. Para la institución académica, se
recomienda conservar el repositorio y la documentación como caso
reutilizable. No se formulan recomendaciones de política pública, porque
el alcance del trabajo no las respalda.

\subsection{Para el desarrollo tecnológico}

Reemplazar las reglas ponderadas por un modelo entrenado recién cuando
haya historial de reservas y reseñas suficiente, manteniendo el diseño
modular que permite hacerlo sin reescribir el sistema. Pasar la
pasarela de pagos (Mercado Pago, hoy en modo de prueba) a producción,
evaluar otros medios de cobro habituales del mercado local y realizar
pruebas de carga una vez que exista un entorno desplegado.

\section{Futuras líneas de investigación}

\textbf{Expansión del estudio actual.} Completar las entrevistas
pendientes, ampliar la cantidad de entrevistados de cada tipo de actor y
la muestra de la encuesta, y observar el uso de la plataforma con
usuarios reales durante un período de varios meses.

\textbf{Nuevas áreas de exploración.} Estudiar desde la economía de
plataformas cómo alcanzar la masa crítica inicial de salones y proveedores
y qué estrategia de precios por lado resulta más conveniente; y, desde la
interacción persona--computadora, si el motivo explicable de cada
recomendación aumenta la confianza de los usuarios.

\textbf{Mejoras metodológicas.} Validar las predicciones del Gemelo
Digital comparándolas con la ejecución de eventos reales y ajustando sus
parámetros; evaluar el recomendador con métricas de precisión definidas de
antemano; y complementar las pruebas automatizadas con pruebas de carga y
de usabilidad.

\textbf{Aplicaciones prácticas y tecnológicas.} Ampliar el alcance a
otros tipos de proveedores y de eventos, a otras ciudades y a una
aplicación móvil; sumar una distribución de mesas sobre un plano a escala
del salón, sugerida en una de las entrevistas, y la simulación de costos
por proveedor.

\section{Cierre}

La organización de eventos se hace hoy con herramientas dispersas, y el
trabajo muestra que es viable concentrarla en una única plataforma. La
plataforma existe, funciona sobre los casos para los que fue probada y
responde a una necesidad que los usuarios relevados reconocen. La
hipótesis se considera respaldada en lo técnico y lo económico, y abierta
en lo que respecta a la mejora medida de eficiencia: lo que queda por
demostrar no es si se puede construir, sino si, usada por personas
reales, efectivamente hace más simple y más eficiente la organización de
un evento. Responder esa pregunta es el paso siguiente natural de este
trabajo.
